\chapter{六極電磁致動器系統建模}
\label{c:modeling}

\section{六極幾何配置}
\label{sec:hexapole_geometry}

% 6 極尖端排列、座標系統定義（Measuring ↔ Actuator）、磁路結構

\section{磁梯度力模型}
\label{sec:force_model}

% 電流基模型 → 磁通基模型推導

\section{過驅動系統}
\label{sec:overactuated}

% 6 輸入 3 輸出、冗餘性、耦合、非線性、位置依賴性

\section{最佳磁通分配}
\label{sec:optimal_allocation}

% 最小化 ‖Φ‖² 的二次規劃、Lagrange 乘子法、LUT 建構

\section{霍爾感測器量測}
\label{sec:hall_sensor}

% 表面安裝 vs 尖端磁通比例關係、準靜態特性

\section{粒子動力學}
\label{sec:particle_dynamics}

% Langevin 方程、Stokes drag、壁效應修正、熱力噪聲
